% Candidate rules: proximity presentation (§4.2.2).  Written by hand; see docs/one-id-survey.md.

\section{Proximity presentation}

\begin{definition}[When proximity presentation is used (C-80)]
  \label{def:c80}
  \uses{thm:photoid-has-portrait}
  Proximity presentation is required for one-to-one face matching when the
  face image was not shared in advance, and should be used whenever
  credentials are shared on demand or offline.

  The touchpoint needs a portrait from the photoid at that moment, so the
  profile depends on the portrait being present.

  Source:
  \href{https://standards.atlassian.net/wiki/spaces/OneID/pages/2947416082/Technical+Implementation+Guide}{One
  ID guide} §4.2.2.
\end{definition}

\begin{definition}[The proximity request (C-84)]
  \label{def:c84}
  \uses{def:proximity-session, def:element}
  The reader must send a device request as ISO/IEC 18013-5 defines it,
  encrypted with the session key. Reader authentication should be present,
  and when it is its structure must follow the same standard.

  A request lists a document type, namespaces and data elements, as in the
  remote profiles.

  External: ISO/IEC 18013-5:2021 §8.3.2.1.2.1, §9.1.4.

  Source:
  \href{https://standards.atlassian.net/wiki/spaces/OneID/pages/2947416082/Technical+Implementation+Guide}{One
  ID guide} §4.2.2, step 3.
\end{definition}

\begin{definition}[The wallet's checks in proximity (C-85)]
  \label{def:c85}
  \uses{def:c84, def:c16}
  \notready
  After decrypting, the wallet must check that the document type matches a
  credential it holds, that the requested data elements match what is stored,
  and must perform reader authentication and validate the session transcript.

  The wallet must authenticate the reader, while a request need only 'should'
  carry reader authentication (C-84).

  Open point: What does a wallet do with a proximity request that has no
  reader authentication: refuse it, or answer it?

  External: ISO/IEC 18013-5:2021 §9.1.4, §9.1.5.

  Source:
  \href{https://standards.atlassian.net/wiki/spaces/OneID/pages/2947416082/Technical+Implementation+Guide}{One
  ID guide} §4.2.2, step 4.
\end{definition}

\begin{definition}[Decrypting the proximity response (C-88)]
  \label{def:c88}
  \uses{def:wallet-respond}
  \notready
  The reader 'should' decrypt the response with the session keys and rebuild
  the session transcript.

  Taken as must: nothing can be validated otherwise.

  Open point: 'Should' for a step every later check needs. Is it must?

  External: ISO/IEC 18013-5:2021 §9.1.5, §9.1.5.1.

  Source:
  \href{https://standards.atlassian.net/wiki/spaces/OneID/pages/2947416082/Technical+Implementation+Guide}{One
  ID guide} §4.2.2, step 5.
\end{definition}

\begin{definition}[Ending the session (C-91)]
  \label{def:c91}
  \uses{def:proximity-session}
  The session shall be terminated when the steps are done or on a condition
  such as a timeout, and all session data made for it must be deleted.

  A session's keys are not used again.

  External: ISO/IEC 18013-5:2021 §9.1.1.4.

  Source:
  \href{https://standards.atlassian.net/wiki/spaces/OneID/pages/2947416082/Technical+Implementation+Guide}{One
  ID guide} §4.2.2, step 7.
\end{definition}

% The proximity rules in Lean.

\begin{definition}[Keys and retrieval in a proximity session (C-82)]
  \label{def:proximity-session}
  \lean{Rules.OneId.Proximity.sessionFindings, Rules.OneId.Proximity.Session, Rules.OneId.Proximity.Retrieval, Rules.OneId.Proximity.SessionFinding, Rules.OneId.Proximity.Session.shared}
  \uses{def:exchange}
  The device makes a key pair for each engagement, the reader makes one before
  each engagement, and data is retrieved from the device and never from a
  server.  The rule looks at the sessions of one deployment together, since a
  reused key shows only across sessions.  The end of a session and the
  deletion of its data are not modelled.

  External: ISO/IEC 18013-5:2021 §8.2.1.1, §9.1.1.4, §9.1.1.5.

  Source: \href{https://standards.atlassian.net/wiki/spaces/OneID/pages/2947416082/Technical+Implementation+Guide}{One ID guide} §4.2.2.
\end{definition}

\begin{theorem}[Accepted proximity sessions use fresh keys and no server]
  \label{thm:proximity-sessions}
  \lean{Rules.OneId.Proximity.accepted_sessions}
  \uses{def:proximity-session}
  In accepted proximity sessions no session retrieves data from a server, no
  reader's key is used twice and no device's engagement key is used twice.
\end{theorem}
\begin{proof}
  \uses{def:proximity-session}
  Server retrieval and a repeated key are each a finding.
\end{proof}
